\documentclass[11pt, a4paper]{article}
\usepackage[utf8]{inputenc}
\usepackage[T1]{fontenc}
\usepackage{graphicx}
\usepackage{hyperref}
\usepackage{geometry}
\usepackage{float}
\usepackage{amsmath, amssymb}
\usepackage{booktabs}
\usepackage{caption}
\usepackage{subcaption}
\usepackage{enumitem}
\usepackage{algorithmic}
\usepackage{algorithm}
\usepackage{xcolor}
\geometry{a4paper, margin=1in}

\hypersetup{
    colorlinks=true,
    linkcolor=blue!70!black,
    citecolor=green!50!black,
    urlcolor=blue!70!black,
}

\title{Homography Synchronization over Graphs:\\
       Iterative, Spectral, and Robust Methods}
\author{}
\date{\today}

\begin{document}

\maketitle

% ===========================================================================
\begin{abstract}
This report describes the design, implementation, and experimental evaluation of
a framework for synchronizing $3\times3$ projective transformations (homographies)
over a graph of pairwise relative measurements.  Two main families of algorithms
are implemented: \emph{iterative} methods based on Madhavan, Fusiello, and
Arrigoni~\cite{madhavan2024}, offering Euclidean, Direction, and Sphere (L1
geodesic mean) averaging with optional IRLS outlier suppression; and
\emph{spectral} (closed-form) methods based on Schroeder, Bartoli, Georgel, and
Navab~\cite{schroeder2011}, providing Locally Scaled Homographies (LSH) and
Globally Scaled Homographies (GSH).  A spanning-tree baseline is included for
reference.  The framework is benchmarked on synthetic graphs---varying graph
size, noise level, edge density, outlier ratio, and topology---as well as on
real images via a full SIFT-based mosaicking pipeline.  Results demonstrate that
Sphere averaging is the most accurate iterative method overall, that IRLS
variants dramatically outperform plain methods in the presence of outliers, and
that the choice of graph topology significantly affects relative algorithm
performance.
\end{abstract}

% ===========================================================================
\section{Introduction}
\label{sec:introduction}

Given a set of $n$ views (images), the \emph{homography synchronization} problem
seeks to recover the unknown absolute projective transformations
$\{X_i\}_{i=1}^{n}$, $X_i \in \mathbb{R}^{3\times3}$, from a set of noisy
pairwise measurements $Z_{ij} \approx X_i \, X_j^{-1}$ available on the edges
of a graph $\mathcal{G} = (\mathcal{V}, \mathcal{E})$.  This problem arises
naturally in image mosaicking, Structure-from-Motion, and multi-view geometry,
wherever relative inter-frame transformations are estimated independently and
must be made globally consistent.

The synchronization solution is inherently defined up to a global right
multiplication: if $\{X_i\}$ is a solution, so is $\{X_i \cdot C\}$ for any
invertible $C$.  All methods in this project operate in $\mathrm{SL}(3)$---the
group of $3\times3$ real matrices with unit determinant---which removes the scale
ambiguity inherent to projective transformations and turns the approximate
consistency relation $Z_{ij} \simeq X_i X_j^{-1}$ into the strict equality
$Z_{ij} = X_i X_j^{-1}$.

% ===========================================================================
\section{Mathematical Formulation}
\label{sec:formulation}

\subsection{The Synchronization Graph}

Let $\mathcal{G} = (\mathcal{V}, \mathcal{E})$ be an undirected graph where
each vertex $i \in \mathcal{V}$ carries an unknown absolute homography
$X_i \in \mathrm{SL}(3)$, and each edge $(i,j) \in \mathcal{E}$ carries a
measured relative homography $Z_{ij}$.  In the noise-free case:
\begin{equation}
    Z_{ij} = X_i \, X_j^{-1}, \qquad \forall\,(i,j) \in \mathcal{E}.
    \label{eq:consistency}
\end{equation}
Both directions are stored: $Z_{ji} = Z_{ij}^{-1}$.

\subsection{SL(3) Normalisation}

Every $3\times3$ matrix $H$ is projected into $\mathrm{SL}(3)$ by dividing by
the cube root of its determinant:
\begin{equation}
    \bar{H} = \frac{H}{\sqrt[3]{\det(H)}}, \qquad \det(\bar{H}) = 1.
\end{equation}
For $3\times3$ real matrices the cube root is always real, making this
normalisation straightforward.

\subsection{Error Metric}

Since the solution is defined up to a global transformation $C$, evaluation
first aligns the estimate to the ground truth by solving $C = \hat{X}_0^{-1}
X_0$ using vertex~$0$ as anchor.  The angular error for vertex $i$ is then:
\begin{equation}
    e_i = \arccos\!\Big(\big|\operatorname{vec}(\hat{X}_i C)^\top
          \operatorname{vec}(X_i)\big|\Big),
    \label{eq:angular_error}
\end{equation}
where $\operatorname{vec}(\cdot)$ denotes unit-norm flattening to $\mathbb{R}^9$.
The absolute value handles the antipodal equivalence on the sphere.  The final
metric is the mean angular error in degrees across all vertices.

% ===========================================================================
\section{Synchronization Algorithms}
\label{sec:algorithms}

\subsection{Iterative Methods (Madhavan et al.~\cite{madhavan2024})}

Each vertex $i$ is updated in turn by computing \emph{neighbour estimates}
\begin{equation}
    X_{i|j} = Z_{ij} \, X_j, \qquad \forall\, j \in \mathcal{N}(i),
    \label{eq:neighbour_est}
\end{equation}
and then averaging these estimates to obtain a new $X_i$.  Vertices are
processed in order of descending degree, since high-degree nodes receive more
constraints and are therefore more stable.  The sweep repeats until the maximum
entry-wise change drops below a tolerance~$\tau$.

Three notions of average are implemented:

\paragraph{Euclidean averaging.}
Normalise each estimate to unit Frobenius norm, compute the weighted
componentwise mean, then re-normalise:
\begin{equation}
    c = \frac{\sum_k w_k \, \bar{h}_k}{\|\sum_k w_k \, \bar{h}_k\|},
\end{equation}
where $\bar{h}_k = \operatorname{vec}(X_{i|j_k}) / \|\operatorname{vec}(X_{i|j_k})\|$.

\paragraph{Direction averaging.}
Finds the eigenvector of
$M = \sum_k w_k \, (h_k h_k^\top) / (h_k^\top h_k)$
corresponding to the largest eigenvalue.  This is L2, scale-invariant.

\paragraph{Sphere averaging (L1 geodesic mean).}
Minimises $\sum_k w_k \arccos(c^\top \bar{h}_k)$ subject to $\|c\|=1$ via the
fixed-point iteration:
\begin{equation}
    c \leftarrow \alpha \sum_k
        \frac{w_k \, \bar{h}_k}{\sqrt{1 - (c^\top \bar{h}_k)^2}},
    \label{eq:sphere_fp}
\end{equation}
iterated until convergence ($\leq 50$ steps, or $\|c_{\text{new}} - c\|_\infty < 10^{-8}$).

\subsubsection{IRLS Outlier Suppression}

When outlier edges are present, an outer \emph{Iteratively Reweighted Least
Squares} (IRLS) loop wraps the inner synchronization sweep.  After each full
convergence, per-edge Cauchy weights are recomputed from the residual:
\begin{equation}
    w_{ij} = \frac{1}{1 + (r_{ij} / s)^2},
    \label{eq:cauchy}
\end{equation}
where $r_{ij}$ is the angular residual (in radians) between $Z_{ij} X_j$ and
$X_i$, and $s$ is the Cauchy scale parameter (typically $0.05$--$0.3$).  Inlier
edges receive weight $\approx 1$, while gross outliers are suppressed toward~$0$.

The inner sweep then restarts with updated weights.  This procedure is repeated
for a fixed number of IRLS iterations (typically 5), requiring no explicit
outlier detection.

\subsection{Spectral Methods (Schroeder et al.~\cite{schroeder2011})}

These methods build a $3n \times 3n$ block matrix from all known inter-frame
homographies and extract the absolute homographies from its null space via SVD.

\paragraph{LSH (Locally Scaled Homographies).}
Builds the block matrix $S$ with entries:
\begin{equation}
    S[k, i] = \frac{\gamma_{i,k}}{\zeta_k} \, H_{i,k},
\end{equation}
where $\gamma_{i,k} = 1$ if the homography from $i$ to $k$ is known,
and $\zeta_k = \sum_i \gamma_{i,k}$ counts known connections to node $k$
(including self).  In the noise-free case $S \cdot U = U$, so $U$ lies in the
null space of $(S - I)$.  The three right singular vectors of $(S-I)$ with the
smallest singular values approximate $U$.

\paragraph{GSH (Globally Scaled Homographies).}
Corrects the weighting bias of LSH by building:
\begin{equation}
    G[k,k] = (1 - \zeta_k)\,I, \qquad G[k,i] = \gamma_{i,k} \, H_{i,k}
    \quad (i \neq k).
\end{equation}
The absolute homographies are again extracted from the null space of $G$ via SVD.

Both methods handle missing data naturally (not all pairs need to be connected).
Their cost is $O((3n)^3)$ in the SVD, making them fast for moderate graph sizes
but less scalable than iterative methods for very large graphs.

\subsection{Spanning-Tree Baseline}

The simplest approach: find a BFS spanning tree rooted at the highest-degree
vertex and chain relative homographies along each path:
\begin{equation}
    X_{\text{child}} = Z_{\text{child},\text{parent}} \cdot X_{\text{parent}}.
\end{equation}
This is $O(n)$ and very fast, but accumulates errors from root to leaves and
does not exploit redundant measurements.

% ===========================================================================
\section{Image Mosaicking Pipeline}
\label{sec:mosaic}

The \texttt{mosaic.py} module implements a full image stitching pipeline that
applies homography synchronization to real images:

\begin{enumerate}[leftmargin=*]
    \item \textbf{Feature extraction:} SIFT keypoints and 128-dimensional
          descriptors are computed on CPU using OpenCV's \texttt{SIFT\_create()}.

    \item \textbf{Descriptor matching:} A GPU-accelerated brute-force matcher
          (\texttt{cv.cuda.DescriptorMatcher}) with L2 norm is used when CUDA is
          available; otherwise a CPU \texttt{BFMatcher} fallback is employed.
          Features are extracted once per image---$O(N)$---rather than
          re-extracted inside the $O(N^2)$ matching loop.

    \item \textbf{Pairwise homography estimation:} Lowe's ratio test (threshold
          $0.7$) filters ambiguous matches from repetitive patterns.  RANSAC
          (inlier tolerance $5.0$~px, minimum $10$ inliers) robustly estimates
          each pairwise homography.

    \item \textbf{Graph construction:} Vertices are images; each successfully
          matched pair becomes an edge carrying its relative homography $Z_{ij}$.

    \item \textbf{Synchronization:} Any of the methods from
          Section~\ref{sec:algorithms} is applied.

    \item \textbf{Mosaic rendering:} Each image is warped into the reference
          frame using the synchronised absolute homographies, via
          GPU-accelerated \texttt{warpPerspective} when available.  A
          translation matrix ensures all coordinates are non-negative.
\end{enumerate}

Quality is evaluated using \emph{mean reprojection error}: for every matched
edge $(i,j)$, inlier points from image $i$ are warped into image $j$ using the
synchronised transform $\hat{Z}_{ji} = \hat{X}_j \hat{X}_i^{-1}$, and the
pixel distance to the actual matched keypoint is measured.  Lower values
indicate better consistency; typical good values are $< 5$~px.

% ===========================================================================
\section{Source Code Overview}
\label{sec:code}

The \texttt{src/} directory contains the following modules:

\begin{description}[leftmargin=*, style=nextline]
    \item[\texttt{graph.py} (1466 lines)]
        Core module implementing the \texttt{Graph} class with adjacency
        structure, SL(3) normalisation, all three iterative averaging functions
        (Euclidean, Direction, Sphere), IRLS weight computation with Cauchy
        robust function, the spectral synchronization methods (LSH, GSH), and
        the spanning-tree baseline.  Also provides utilities for synthetic graph
        construction (random Erd\H{o}s--R\'{e}nyi, linear chain, circular ring,
        2-D grid, and multi-lane topologies), ground-truth generation in SL(3),
        noise injection, outlier edge insertion, angular error computation, and
        graph visualization via NetworkX and Matplotlib.

    \item[\texttt{mosaic.py} (601 lines)]
        Full image mosaicking pipeline: SIFT feature extraction, GPU/CPU
        BFMatcher descriptor matching, Lowe's ratio test, RANSAC homography
        estimation, graph construction from images, reprojection error
        computation, mosaic rendering with GPU-accelerated warp, and
        command-line interface with configurable synchronization method,
        averaging variant, IRLS iterations, and Cauchy scale.

    \item[\texttt{benchmark.py} (831 lines)]
        Main benchmarking script running five experiments (Sections
        \ref{sec:exp1}--\ref{sec:exp5}).  Implements the method registry
        (6~plain methods $+$ 3~IRLS variants), a single-trial runner
        supporting all topologies, experiment sweep functions, consistent
        plotting with per-method visual styles, a combined results tiler,
        and a command-line \texttt{--from N} option for resuming from a
        specific experiment.

    \item[\texttt{benchmark\_noise\_topology.py} (549 lines)]
        Joint noise $\times$ topology benchmark (Section~\ref{sec:exp_topo}).
        Sweeps noise levels across six graph topologies (linear bw=1, linear
        bw=3, circular chord=3, grid 8-connected, multilane 3~lanes, random
        E-R) and produces a grid of error-vs-noise line plots, a per-method
        topology-comparison plot, and a winner heatmap.

    \item[\texttt{compress.py} (49 lines)]
        Utility for resizing and JPEG-compressing a dataset of images to
        accelerate processing and reduce memory usage during mosaicking
        (default: $25\%$ scale, quality~100).
\end{description}

% ===========================================================================
\section{Experimental Setup}
\label{sec:setup}

All synthetic experiments use 50 independent trials per data point, and the
reported metric is the \emph{median} angular error (in degrees) to reduce
sensitivity to rare divergence events.  Unless otherwise noted, the default
graph topology is random (Erd\H{o}s--R\'{e}nyi with edge dropout), with $n=25$
nodes.

Six synchronization methods are compared in Experiments 1--3:
\textbf{Tree} (baseline), \textbf{Sphere}, \textbf{Euclidean},
\textbf{Direction} (iterative), \textbf{LSH}, and \textbf{GSH} (spectral).
Experiment~4 adds three IRLS variants: \textbf{Sphere-IRLS},
\textbf{Euclidean-IRLS}, and \textbf{Direction-IRLS} (5~IRLS iterations,
Cauchy scale $s = 0.15$).

% ===========================================================================
\section{Experimental Results}
\label{sec:results}

% ------- Experiment 1 -------
\subsection{Experiment 1: Varying Graph Size}
\label{sec:exp1}

\textbf{Setup:} $n \in \{10, 20, 30, 50, 75, 100\}$, $\sigma = 0.05$,
$\rho = 0.5$ (hole density), $\gamma = 0$ (no outliers).

\begin{figure}[H]
    \centering
    \includegraphics[width=\textwidth]{results/benchmark/exp1_vary_nodes.png}
    \caption{Experiment 1: Angular error (left) and execution time (right)
             as a function of graph size.}
    \label{fig:exp1}
\end{figure}

\textbf{Observations:}  As the graph grows, more redundant measurements become
available, and all iterative/spectral methods improve---Sphere averaging
consistently achieves the lowest error (from ${\sim}8^\circ$ at $n{=}10$ down to
${\sim}2^\circ$ at $n{=}100$).  The spanning-tree baseline \emph{worsens} with
size (from ${\sim}13^\circ$ to ${\sim}16^\circ$), confirming its error
accumulation from root to leaves.  On execution time, LSH/GSH remain the
fastest synchronization methods (single SVD solve), while Sphere is the slowest
due to its per-node fixed-point iteration; the tree baseline is near-instant.

% ------- Experiment 2 -------
\subsection{Experiment 2: Varying Noise Level}
\label{sec:exp2}

\textbf{Setup:} $\sigma \in \{0.01, 0.02, 0.04, 0.06, 0.08, 0.10, 0.12,
0.15\}$, $n = 25$, $\rho = 0.5$, $\gamma = 0$.

\begin{figure}[H]
    \centering
    \includegraphics[width=\textwidth]{results/benchmark/exp2_vary_noise.png}
    \caption{Experiment 2: Angular error (left) and execution time (right)
             as noise increases.}
    \label{fig:exp2}
\end{figure}

\textbf{Observations:}  All methods degrade with increasing noise, but the
iterative and spectral approaches degrade more gracefully than the spanning tree.
This is because iterative and spectral methods exploit all available edges to
average out independent noise contributions, whereas the spanning tree propagates
each noisy measurement directly along the path from root to leaf with no
cross-checking.  Sphere averaging maintains the lowest error throughout (from
${\sim}1^\circ$ at $\sigma{=}0.01$ to ${\sim}16^\circ$ at $\sigma{=}0.15$).
LSH/GSH degrade more rapidly at high noise and converge toward the tree
baseline: the spectral block-matrix solve is a global weighted least-squares
fit that gives equal weight to every measurement, so large noise corrupts the
system uniformly and the SVD solution loses accuracy without any robust
averaging to compensate.  Execution times are largely independent of noise
because convergence speed is governed by graph structure (number of neighbours,
degree distribution), not by the magnitude of the perturbations.

% ------- Experiment 3 -------
\subsection{Experiment 3: Varying Hole Density}
\label{sec:exp3}

\textbf{Setup:} $\rho \in \{0.0, 0.1, 0.2, 0.3, 0.5, 0.7, 0.8, 0.9\}$,
$n = 25$, $\sigma = 0.05$, $\gamma = 0$.

\begin{figure}[H]
    \centering
    \includegraphics[width=\textwidth]{results/benchmark/exp3_vary_holes.png}
    \caption{Experiment 3: Angular error (left) and execution time (right)
             as the fraction of missing edges increases.}
    \label{fig:exp3}
\end{figure}

\textbf{Observations:}  With denser graphs ($\rho \to 0$), all methods benefit
from abundant redundancy.  As more edges are removed, less redundancy is
available for error compensation and all methods converge toward spanning-tree
performance; in the extreme limit (minimum spanning tree), every method
necessarily reduces to a single chain of chained measurements with no
cross-checking.  Sphere maintains the best accuracy across the full range
because the L1 geodesic mean down-weights the influence of any single noisy
measurement, so even a sparse set of neighbours is exploited efficiently.
The spectral methods (LSH/GSH) handle missing data natively and remain
competitive until very high sparsity ($\rho > 0.7$): absent edges simply
contribute zero blocks to the block matrix, leaving the SVD problem
well-posed without requiring any special treatment.  Iterative execution time
decreases with sparsity (fewer neighbours to average), while spectral time is
roughly constant (dominated by the fixed $3n \times 3n$ SVD regardless of
how many blocks are zero).


% ------- Experiment 4 -------
\subsection{Experiment 4: Varying Outlier Density}
\label{sec:exp4}

\textbf{Setup:} $\gamma \in \{0.0, 0.05, 0.1, 0.2, 0.3, 0.4, 0.5, 0.6\}$,
$n = 25$, $\sigma = 0.05$, $\rho = 0.2$.

\begin{figure}[H]
    \centering
    \includegraphics[width=\textwidth]{results/benchmark/exp4_vary_outliers.png}
    \caption{Experiment 4: Plain methods (solid) vs.\ IRLS variants (dotted).
             Left: angular error; right: execution time.}
    \label{fig:exp4}
\end{figure}

\textbf{Observations:}  This experiment clearly demonstrates the value of IRLS
robust reweighting.  All plain methods degrade rapidly beyond $\gamma{=}0.1$,
reaching ${\sim}60$--$70^\circ$ error at $\gamma{=}0.6$---effectively random.
This collapse occurs because plain averaging weights every neighbor equally:
even a small fraction of outlier edges (random homographies) pulls the mean
far from the true estimate, and without downweighting these corrupt
measurements the iteration converges to a wrong fixed point.  In contrast,
the \textbf{Sphere-IRLS} variant maintains ${\sim}5$--$8^\circ$ error even
at $\gamma{=}0.5$, only degrading significantly at $\gamma{=}0.6$: Cauchy
reweighting progressively suppresses edges whose residuals are large (i.e.,
inconsistent with the current estimate), allowing the method to identify and
ignore outliers implicitly.  Euclidean-IRLS is the second-best robust method,
slightly behind Sphere-IRLS because the L2 mean is more sensitive to residual
outlier pull than the L1 sphere mean.  The spectral methods (LSH/GSH) are not
inherently robust and deteriorate as rapidly as plain iterative methods because
the SVD null-space extraction is a global least-squares problem that cannot
localize or down-weight individual corrupt edges.  The IRLS cost is modest:
Sphere-IRLS takes ${\sim}4$~s at $\gamma{=}0.6$ compared to ${\sim}0.5$~s
for plain Sphere.


% ------- Experiment 5 -------
\subsection{Experiment 5: Real Image Data}
\label{sec:exp5}

\textbf{Setup:} 10 images from a compressed dataset, evaluated by mean
reprojection error (pixels).

\begin{figure}[H]
    \centering
    \includegraphics[width=\textwidth]{results/benchmark/exp5_real_data.png}
    \caption{Experiment 5: Reprojection error (left) and execution time
             (right) on real images.}
    \label{fig:exp5}
\end{figure}

\textbf{Observations:}  On real data, the \textbf{Tree} baseline achieves the
lowest reprojection error (${\sim}2.2$~px), while the iterative and spectral
methods cluster around $2.8$--$3.4$~px.  This counterintuitive result has a
clear cause: the real dataset consists of sequentially captured images with
predominantly linear overlap, and the pairwise homographies estimated by
SIFT+RANSAC are already clean (no gross outliers, low reprojection error per
pair).  In this setting the spanning tree is the globally optimal solution
because the graph has minimal cycle redundancy and no corrupt edges to
average out.  The iterative and spectral methods are designed to exploit
cycle redundancy to correct errors; when that redundancy is absent and the
measurements are clean, they instead accumulate the effect of small
inter-frame inconsistencies over multiple rounds of global optimization,
producing a slightly worse result than the simple chain.  Sphere averaging
is the slowest (${\sim}0.16$~s) while LSH and GSH are near-instantaneous.


% ------- Topology Experiments -------
\subsection{Experiment 6: Joint Noise $\times$ Topology Analysis}
\label{sec:exp_topo}

\textbf{Setup:}  $\sigma \in \{0.02, 0.04, 0.10, 0.15, 0.20, 0.25\}$, $n=25$,
$\rho=0.3$, 20~trials, across six topologies: Linear (bw=1), Linear (bw=3),
Circular (chord=3), Grid (8-connected), Multilane (3~lanes), Random (E-R).

\begin{figure}[H]
    \centering
    \includegraphics[width=\textwidth]{results/benchmark_noise_topology/exp_noise_topology_grid.png}
    \caption{Error vs.\ noise across six graph topologies.  Each subplot shows
             all methods on one topology; lower is better.}
    \label{fig:topo_grid}
\end{figure}

\begin{figure}[H]
    \centering
    \includegraphics[width=0.95\textwidth]{results/benchmark_noise_topology/exp_noise_topology_heatmap.png}
    \caption{Winner heatmap: best method at each (topology, noise) cell.
             Cell colour encodes best achievable error; text shows the
             winning method abbreviation and error.}
    \label{fig:topo_heatmap}
\end{figure}

\textbf{Observations from the grid plot (Figure~\ref{fig:topo_grid}):}
\begin{itemize}[leftmargin=*]
    \item \textbf{Linear (bw=1):}  All methods perform poorly because the
          pure chain provides no cycle redundancy---there is no alternative
          path between any two nodes, so no cross-checking of measurements
          is possible and errors accumulate exactly as in the spanning tree.
          Tree and GSH are competitive at low noise; all methods saturate near
          $60$--$65^\circ$ at high noise.

    \item \textbf{Linear (bw=3):}  The wider bandwidth adds shortcut edges
          that create short cycles, enabling limited cross-checking and
          improving all methods.  Tree still wins at low noise because the
          extra edges offer little benefit when individual measurements are
          already accurate; at moderate noise Sphere and Direction recover as
          they can now average over genuinely redundant information.

    \item \textbf{Circular (chord=3):}  Loop closure introduces non-trivial
          cycles that provide global consistency constraints (any closed-loop
          product of relative homographies must equal identity), giving all
          methods additional signal to correct drift.  LSH and Euclidean
          excel at low noise because the linear system is already
          well-conditioned; Direction and Sphere perform better at moderate
          noise where robustness to individual measurements matters more.

    \item \textbf{Grid (8-connected):}  Dense local redundancy (every node
          has up to 8 neighbours) benefits both spectral and iterative methods
          equally, since each node receives many independent constraint
          equations from different directions.  Euclidean excels at low noise
          (many measurements with small error → accurate L2 mean); Sphere at
          moderate noise (L1 geodesic mean down-weights the noisier
          neighbours).

    \item \textbf{Multilane (3~lanes):}  Sphere dominates across all noise
          levels because the multilane structure provides high local
          connectivity \emph{and} long-range cross-lane edges, creating rich
          cycle structure.  This is exactly the regime where the L1 geodesic
          mean is most effective: many redundant neighbours with varying
          quality, where robustly down-weighting the worst ones yields a large
          accuracy gain over the L2 mean.

    \item \textbf{Random (E-R):}  With high and uniform average connectivity,
          every node has many neighbours from which to form an accurate
          estimate.  Sphere achieves the best performance overall because the
          random degree distribution means some neighbours are more reliable
          than others, and the L1 mean naturally suppresses their influence.
\end{itemize}


\textbf{Observations from the heatmap (Figure~\ref{fig:topo_heatmap}):}
No single method dominates all conditions.  \textbf{Sphere} wins most cells on
random, multilane, and grid topologies.  \textbf{Tree} is competitive on linear
graphs, especially at moderate-to-high noise where iterative methods suffer from
the lack of redundancy.  Spectral methods (GSH, LSH) win only on sparse
topologies at low noise.  \textbf{Euclidean} is strong on grid and circular
topologies at low noise.
% ------- IRLS Norm Experiment -------
\subsection{Attempted Experiment 1: IRLS Norm variation}
\label{sec:exp_IRLS_Norm}
As an attempt to improve the Sphere method,  we implemented the \textbf{L1 norm} (using Weiszfeld's algorithm within the IRLS framework) to improve the robustness of the iterative synchronization method. Unfortunately, the experimental results did not yield the expected improvements in accuracy compared to the baseline Cauchy loss. 
\begin{figure}[H]
    \centering
    \includegraphics[width=\textwidth]{results/benchmark L_1/l1_exp2_noise.png}
    
    \vspace{0.5cm} % Piccolo spazio tra le due immagini
    
    \includegraphics[width=\textwidth]{results/benchmark L_1/l1_exp1_outliers.png}
    
    \caption{Top: Performance at varying noise levels. Bottom: Performance at varying outlier densities.}
    \label{fig:benchmark_l1}
\end{figure}

Comparing the iterative method with Cauchy loss versus the L1 norm, we observed two main trends:
\begin{itemize}
    \item \textbf{Accuracy:} At typical noise levels ($\sigma = 0.05$), the L1 norm produces a significantly higher angular error than Cauchy ($\sim 8^\circ$ vs. $\sim 5^\circ$).
    \item \textbf{Speed:} The L1 norm converges systematically faster, keeping execution times consistently lower.
\end{itemize}

\textbf{Why it failed.}
The higher error of the L1 norm is due to its \textbf{statistical efficiency} regarding the underlying noise model. 
In our framework, measurement noise is modeled as \textbf{Gaussian}. According to the Gauss-Markov theorem, the optimal estimator for Gaussian noise is the \textbf{L2 Norm}.
\begin{itemize}
    \item \textbf{Cauchy Loss:} For inliers with small residuals ($r \ll c$), the Cauchy IRLS weights $w = 1 / (1 + (r / c)^2)$ approach 1. Consequently, Cauchy asymptotically \textit{behaves like L2} for the inliers, achieving optimal efficiency against Gaussian noise.
    \item \textbf{L1 Norm:} The L1 IRLS weights diverge near zero ($w = 1/|r|$), forcing the solution towards the \textit{geometric median}. While excellent at rejecting gross outliers, the median is a sub-optimal estimator (with higher variance) for Gaussian distributions, resulting in a higher residual error when dealing with noisy inliers.
\end{itemize}
Despite being less accurate, L1 is computationally faster. This is because its diverging weights ($1/|r|$) act as an aggressive \textit{hard-attention} mechanism. The algorithm rapidly locks onto a subset of edges, driving their residuals to zero and quickly freezing the graph's geometry. Cauchy, on the other hand, seeks a "soft" and continuous equilibrium among all inliers, naturally requiring more iterations (sweeps) to fully converge.
% ------- Lie Branch -------
\subsection{Attempted Experiment 2: Lie-Algebraic (Riemannian) Synchronization}
\label{sec:exp_lie}

As part of the project, a Riemannian synchronization method on $\mathrm{SL}(3)$
was implemented in a separate \texttt{lie} branch alongside the corresponding
benchmark script \texttt{benchmark\_sphere\_vs\_lie.py}.  The idea is to replace
the extrinsic sphere-based averaging---which treats each homography as a
9-dimensional unit vector in $\mathbb{R}^9$---with an \emph{intrinsic} update
that respects the curved geometry of the Lie group:
\begin{equation}
    X_i^{\text{new}} = X_i \cdot \exp\!\Bigl(
        \tfrac{1}{\sum_j w_j} \sum_{j \in \mathcal{N}(i)}
        w_j \, \log\!\bigl(X_i^{-1} Z_{ij} X_j\bigr)
    \Bigr),
\end{equation}
where $\log(\cdot)$ and $\exp(\cdot)$ denote the matrix logarithm and exponential
on $\mathrm{SL}(3)$, and $\mathfrak{sl}(3)$ (the space of $3\times3$ traceless
matrices) serves as the Lie algebra in which the averaging is performed.

\textbf{Why it failed.}
The fundamental obstacle is the matrix logarithm.  \texttt{scipy.linalg.logm}
computes the \emph{principal} matrix logarithm via Schur decomposition, which is
only real-valued for matrices whose eigenvalues have no negative real part.
Generic $\mathrm{SL}(3)$ matrices---such as the randomly generated ground-truth
homographies used in the benchmarks---frequently have negative real eigenvalues,
causing \texttt{logm} to return complex values and making the Lie-algebra update
undefined.  A more robust fallback (\texttt{\_sl3\_log\_robust}) was developed,
which iteratively takes matrix square roots to bring the input into a small
neighborhood of the identity before applying \texttt{logm} (following the
$H \approx (H^{1/2^k})^{2^k}$ identity).  However, this approach also fails
when any intermediate square root lands on a matrix with negative real eigenvalues,
since real square roots do not exist in that case.  In practice, a significant
fraction of neighbor estimates $X_i^{-1} Z_{ij} X_j$ fell outside the domain of
the real logarithm, forcing those neighbors to be skipped and degrading the update
quality.  Combined with the $3$--$5\times$ per-sweep computational overhead of
\texttt{logm} and \texttt{expm} relative to the sphere fixed-point iteration, the
Lie method did not achieve better accuracy than Sphere averaging and was therefore
excluded from the main experiments.

% ===========================================================================
\section{Real-Life Results}
\label{sec:real_life}

The mosaic of a building in Piacenza constructed using the Sphere averaging method is shown below.

\begin{figure}[H]
    \centering
    \includegraphics[width=0.8\textwidth]{results/real_life_results/real_life_homography.png}
    \caption{Mosaic of a building in Piacenza, constructed using the Sphere averaging method.}
    \label{fig:real_life_mosaic}
\end{figure}

% ===========================================================================
\section{Summary of Findings}
\label{sec:summary}

\begin{table}[H]
\centering
\caption{Comparative summary of synchronization methods.}
\label{tab:summary}
\small
\begin{tabular}{@{}lp{5.5cm}p{5.5cm}@{}}
\toprule
\textbf{Family} & \textbf{Strengths} & \textbf{Weaknesses} \\
\midrule
\textbf{Iterative} (Sphere) &
    Best accuracy overall; inherent L1 robustness; handles large graphs
    efficiently ($O(n \cdot d)$ per iteration) &
    Requires multiple iterations to converge; slowest per-node computation
    due to fixed-point loop \\
\addlinespace
\textbf{Iterative} (Euclidean) &
    Competitive accuracy; simpler than Sphere; strong on grid/circular
    topologies &
    L2 averaging; less robust than Sphere \\
\addlinespace
\textbf{Iterative} (Direction) &
    Scale-invariant; good on dense graphs &
    L2-based eigenvector; less robust to outliers \\
\addlinespace
\textbf{IRLS variants} &
    Dramatic improvement under outliers ($\gamma \leq 0.5$); Cauchy
    weighting suppresses gross errors without explicit detection &
    $5\times$ overhead from repeated convergence cycles \\
\addlinespace
\textbf{Spectral} (LSH/GSH) &
    Single SVD solve---no iterations; handles missing data natively;
    fastest synchronization step &
    $O((3n)^3)$ SVD cost; not inherently robust to outliers; less
    accurate than iterative on dense graphs \\
\addlinespace
\textbf{Spanning Tree} &
    Fastest method (single BFS, $O(n)$); works well on low-noise linear
    data &
    Accumulates errors from root to leaves; does not exploit redundant
    measurements \\
\addlinespace
\textbf{Lie-algebraic} (Riemannian) &
    Intrinsic updates respect the curved geometry of $\mathrm{SL}(3)$;
    theoretically faster convergence for large homographies &
    Principal matrix logarithm (\texttt{logm}) is undefined over the reals
    for matrices with negative eigenvalues, which are common in $\mathrm{SL}(3)$; \\
\bottomrule
\end{tabular}
\end{table}

\noindent\textbf{Key takeaways:}
\begin{enumerate}[leftmargin=*]
    \item \textbf{Sphere averaging is the best general-purpose method}, with the
          lowest error on random and multilane topologies across all noise levels.
    \item \textbf{IRLS is essential for outlier-contaminated data}: Sphere-IRLS
          maintains $<10^\circ$ error even with 50\% outlier edges.
    \item \textbf{Topology matters}: linear chains with low bandwidth severely
          limit all methods due to minimal redundancy; adding loop closure
          (circular) or cross-connections (multilane, grid) dramatically improves
          performance.
    \item \textbf{Spectral methods are efficient but fragile}: LSH/GSH excel on
          clean, moderately sparse graphs but degrade rapidly with noise or
          outliers.
    \item \textbf{The spanning tree can be surprisingly competitive} on
          sequential (linear) real-world datasets with low noise.
    \item \textbf{Intrinsic Lie-algebraic averaging is impractical for general
          $\mathrm{SL}(3)$}: the principal matrix logarithm fails on matrices
          with negative real eigenvalues, which occur frequently in this group,
          making the method unreliable without a more sophisticated (and
          expensive) geodesic computation.
\end{enumerate}

% ===========================================================================

\begin{thebibliography}{9}

\bibitem{madhavan2024}
R.~Madhavan, A.~Fusiello, and F.~Arrigoni,
``Synchronization of Projective Transformations,''
2024.

\bibitem{schroeder2011}
P.~Schroeder, A.~Bartoli, P.~Georgel, and N.~Navab,
``Closed-Form Solutions to Multiple-View Homography Estimation,''
in \emph{Proc.\ IEEE Workshop on Applications of Computer Vision (WACV)}, 2011.

\end{thebibliography}

\end{document}
